\section{Introduction}\label{sec:introduction}


\subsection{Motivation}\label{sec:motivation}

Random values have been a fundamental building block of Web2 applications, powering everything from gaming mechanics to cryptographic protocols. As we transition to Web3, we must maintain these essential capabilities while upholding our core principles of decentralization and security. However, blockchain's inherently deterministic and public nature presents a unique challenge in preserving the unpredictability of random values.

This fundamental limitation creates what we term the ``\textbf{Onchain Randomness Problem}'': as soon as entropy becomes available on the blockchain, all outcomes dependent on that entropy become predictable by any observer. Traditional attempts to generate randomness within blockchain systems have led to various compromises in security, integrity, or decentralization, undermining the very principles Web3 aims to uphold.

\subsection{Bringing Back \texttt{math.random}}\label{sec:math_random}

On AO the problem is now concrete. Release v0.11-RC7 of HyperBEAM, the node software that runs AO processes, replaces \texttt{math.random} in its Lua runtime with a sandboxed placeholder, together with \texttt{os.time}, \texttt{os.clock} and \texttt{os.date}. Each of these returns a value that depends on the node running it: its random seed, its clock, its time zone. That makes them unsafe in two ways:

\begin{itemize}
    \item \textbf{Forks}: an AO process is defined by replaying its messages. If two nodes, or the same node on a replay, draw different numbers, they compute different states for the same process, and its history forks.
    \item \textbf{Manipulation}: whoever controls the seed controls the draw. A node operator can choose or observe the seed and predict or steer every roll, which makes node-local randomness worthless for anything of value.
\end{itemize}

RandAO brings \texttt{math.random} back safely. A single call, \texttt{rng.install(msg)}, gives \texttt{math.random} a deterministic stream derived from the message, identical on every node and every replay, so existing code keeps working at no cost. When a draw carries value, \texttt{rng.request} upgrades the same code to verifiable randomness from the network, at the security level the application chooses.

\subsection{Our Contribution}\label{sec:our_contribution}

RandAO is a suite of cryptographic tools for the new age of the internet, delivered through one library and one HyperBEAM device. Its first tool is tiered randomness; time-lock encryption follows. The flagship is the second version of the time-lock protocol RandAO was founded on, rebuilt for HyperBEAM:

\begin{itemize}
    \item \textbf{Any-Honest Security}: The output is unpredictable and unbiasable as long as at least one committed validator is honest.

    \item \textbf{Withholding Cannot Bias}: Once commitments close, every contribution is fixed. A validator that refuses to reveal is cracked; the worst it can do is delay the round.

    \item \textbf{Verification on the Node}: A reveal carries the puzzle's one-time trapdoor, so the node checks it with one modular exponentiation through the \texttt{timelock@1.0} device. Cracks carry a Wesolowski proof. No verifier fleet is needed.

    \item \textbf{Self-Contained}: Validators generate their own puzzles and hold their own keys. There is no shared secret, no key ceremony and no dependence on outside randomness.

    \item \textbf{Any Security Level, Per Call}: Applications choose between free local randomness, a fast threshold-signature beacon and the time-lock beacon on every call. Moving up a tier changes one option, not the code.
\end{itemize}
